%%%%%%%%%%%%%%%%%
% Визуальная версия CV (двухколоночная, altacv.cls v1.1.5).
% Компилятор: pdfLaTeX.
% Требуется файл altacv.cls из этого же проекта Overleaf.
% Источник фактов: resume-facts.yaml (подтверждено 2026-10-23, обновлено 2026-10-24).
%%%%%%%%%%%%%%%%%

\documentclass[9pt,a4paper,ragged2e]{altacv}

\geometry{left=1.1cm,right=1.1cm,top=1.2cm,bottom=1.2cm,columnsep=1cm}
\usepackage[english,russian]{babel}
\usepackage{hyperref}
\usepackage{paracol}

\ifxetexorluatex
  \setmainfont{Carlito}
\else
  \usepackage[utf8]{inputenc}
  \usepackage[T1]{fontenc}
  \usepackage[default]{lato}
\fi

\definecolor{VividPurple}{HTML}{000000}
\definecolor{SlateGrey}{HTML}{2E2E2E}
\definecolor{LightGrey}{HTML}{2E2E2E}
\colorlet{heading}{VividPurple}
\colorlet{accent}{VividPurple}
\colorlet{emphasis}{SlateGrey}
\colorlet{body}{LightGrey}

\renewcommand{\itemmarker}{{\small\textbullet}}
\renewcommand{\ratingmarker}{\faCircle}

\begin{document}
\name{Ачикян Эдуард}
\personalinfo{%
  \email{\href{mailto:ed9119@mail.ru}{ed9119@mail.ru}}
  \printinfo{\faTelegram}{\href{https://t.me/edooam}{@edooam}}
  \phone{+7-926-957-73-83}
  \location{Москва}
  \github{\href{https://github.com/EdA20}{github.com/EdA20}}
}

\begin{fullwidth}
\makecvheader
\end{fullwidth}

\AtBeginEnvironment{itemize}{\medium}

\columnratio{0.5}
\begin{paracol}{2}

% ==================== ЛЕВАЯ КОЛОНКА ====================
\cvsection{Опыт}

\cvevent{Сбер, Аналитическое управление}{Аналитик индексной команды}{Май 2024 -- настоящее время}{Москва}
\skipsmall
\begin{itemize}[leftmargin=12pt]
  \item Разработка инвестиционных индексов Сбера - написание методологии расчёта, реализация алгоритма расчета на Python, согласование с другими департаментами (правовой, комплаенс, риски, трейдинг), публикация значений на портале
  \item Передача сигналов ребалансировки и хеджированию позиций для трейдинг деска по индексам
  \item Разработка на python отчета по контролю за опубликованными значениями
  \item Сбор рыночных данных через REST API (Interfax, MOEX)
  \item С использованием LLM переработал Python-библиотеку для работы с API внутренней системы
\end{itemize}

\divider

\cvevent{Райффайзенбанк}{Аналитик рыночных рисков}{Июнь 2023 -- Май 2024}{Москва}
\skipsmall
\begin{itemize}[leftmargin=12pt]
  \item Построение структуры ПИФ и ДУ в разрезе активов по каждому клиенту
  \item Контроль состава портфелей ДУ с учётом риск-профиля клиентов
  \item Автоматизация сбора рыночных данных через REST API Interfax, Cbonds, MOEX (Airflow, Gitlab)
  \item Интеграция новых таблиц из БД (SQL) в отчеты, взаимодействие с IT подразделением как бизнес-заказчик (Jira, Confluence)
  \item Автоматизация построения отчетности по рискам (VAR, лимит контрагентского риска)
\end{itemize}

\divider

\cvevent{Банк Открытие}{Специалист по рыночным рискам}{Июнь 2022 -- Июнь 2023}{Москва}
\skipsmall
\begin{itemize}[leftmargin=12pt]
  \item Ускорение расчёта ТСС структурных продуктов методом Монте-Карло, реализовав алгоритм на Python: 15 минут вместо 5 часов
  \item Расчет риск-сценариев и греков по структурным продуктам (delta, gamma vega)
  \item Расчет VaR по портфелю облигаций через G-spread
  \item Сбор рыночных данных (selenium + requests), построение базы данных (SQL)
\end{itemize}

\switchcolumn

% ==================== ПРАВАЯ КОЛОНКА ====================

\cvsection{Навыки}
\begin{itemize}[leftmargin=12pt]
  \item \textbf{Python:} pandas, numpy, scikit-learn, pytorch, matplotlib, requests, selenium, langgraph/langchain
  \item \textbf{Risk:} VaR, CVaR, Stress Test, Monte-Carlo, Sharpe/Sortino ratio
  \item \textbf{Данные:} SQL, Airflow, REST API, Hive, Hadoop, PySpark, Greenplum, ClickHouse
  \item \textbf{Инфраструктура:} Docker, FastAPI, Git/GitLab, DVC, MLflow, Hydra, pre-commit, Jira, Confluence
  \item \textbf{Языки:} русский; английский --- B2 (Upper-Intermediate)
\end{itemize}

\cvsection{Проекты}

\cvevent{Portfolio Risk Agent}{LLM-агент для риск-анализа портфеля}{2026}{}
AI-агент на LangGraph, который оценивает качество ребалансировки портфеля с точки зрения рыночного риска и новостного контекста (\hyperlink{https://github.com/EdA20/PorfolioRiskAgent}{ссылка})

\divider

\cvevent{Investment Portfolio Construction}{Магистерская ВКР НИУ ВШЭ}{2025}{}
ML-стратегия, аллоцирующая капитал между индексом IMOEX и денежным рынком (\hyperlink{https://github.com/EdA20/Investment-Portfolio-Construction}{ссылка})

\cvsection{Образование}

\cvevent{НИУ ВШЭ}{Магистратура, «Нейронные сети»}{Сентябрь 2022 -- Июнь 2025}{}
\begin{itemize}[leftmargin=12pt]
    \item \textbf{ML/DL/RL:} sklearn, pytorch, gymnasium
    \item \textbf{Big Data:} Hive, Hadoop, PySpark, Greenplum, ClickHouse
    \item \textbf{MLOps:} Git, DVC, pre-commit, Hydra, MLflow, wandb
    \item \textbf{DevOps:} Docker, FastAPI, Airflow
\end{itemize}

\divider

\cvevent{НИУ ВШЭ}{Бакалавриат, «Экономика»}{Сентябрь 2018 -- Июнь 2022}{}
\begin{itemize}[leftmargin=12pt]
    \item Теория вероятностей и математическая статистика
    \item Эконометрика
    \item Финансовый рынки
    \item Корпоративные финансы
\end{itemize}

\divider

\cvevent{МФТИ}{Курс «Глубокое обучение и компьютерное зрение»}{Сентябрь 2025 -- Февраль 2026}{}

\end{paracol}
\end{document}
